\section{Introduction}
\label{sec:intro}

Humans reason through analogy constantly.
We call sensors a ``hat,'' describe organizational structure with body metaphors (``head,'' ``hands,'' ``heart''), and map code repositories to trees with ``branches'' and ``trunks'' \citep{lakoff1980metaphors, hofstadter2001analogy}.
These analogies are \emph{partial isomorphisms}: structure-preserving maps between domains that hold for some properties but not all \citep{gentner1983structure}.
Large language models, trained to predict text through compression of vast corpora, face the same pressure to find useful abstractions across domains---but they are not constrained to discover the same mappings that humans use.

{\bf Do LLMs discover their own, counterintuitive partial isomorphisms?}
The superposition literature establishes that neural networks routinely compress more features than they have dimensions, entangling representations of distinct concepts \citep{elhage2022superposition, scherlis2022polysemanticity}.
The OthelloGPT work provides a striking example: the model encodes board state as MINE/YOURS rather than BLACK/WHITE, conflating two human categories into one internal concept \citep{nanda2023emergent}.
However, no prior work has systematically catalogued which unexpected concept pairings LLMs create, measured how surprising these pairings are relative to human intuition, or tested whether they confer reasoning advantages.

We address this gap with a three-pronged investigation.
We define a \emph{surprise score} that quantifies the discrepancy between LLM embedding similarity and human conceptual similarity for concept pairs across 12 domains.
Using this metric, we identify 162 significantly surprising cross-domain pairs out of 6,600 tested---pairs like \emph{ingredient}$\leftrightarrow$\emph{enzyme} (both inputs to a transformation), \emph{foundation}$\leftrightarrow$\emph{anchor} (both provide stability), and \emph{stack}$\leftrightarrow$\emph{layer} (both involve ordered accumulation).
We then validate these mappings through two complementary lenses: behavioral probing shows that unexpected analogies improve \gptfour's reasoning quality by 20.7\% over no-hint baselines, and mechanistic analysis reveals that surprising pairs share significantly more MLP activation overlap in \pythia (Cohen's $d = 0.81$, $p = 0.013$), concentrated in semantic processing layers 5--17.

Our contributions are:
\begin{itemize}[leftmargin=*,itemsep=0pt,topsep=0pt]
    \item We propose the \surprisescore metric for systematically discovering unexpected conceptual mappings in LLM embedding spaces, identifying 162 significantly surprising cross-domain concept pairs.
    \item We demonstrate that these unexpected analogies are not structurally empty: they improve reasoning quality when used as hints and are validated as structurally sound by \gptfour.
    \item We provide mechanistic evidence that surprising concept pairs share activation patterns in \pythia, with effects concentrated in semantic processing layers, linking behavioral observations to internal representations.
\end{itemize}

We describe related work in \secref{sec:related}, our methodology in \secref{sec:method}, results in \secref{sec:results}, and discuss implications and limitations in \secref{sec:discussion}.
